\documentclass{article}
\usepackage{graphicx} % Required for inserting images
\usepackage{amsmath} % Equations for Part 3

\title{The effect of State Sponsorship on the Lethality of Terrorist Organizations, Measured by the Number of Civilian Fatalities Attributed to Each Organization from 1998 to 2005.}
\author{Katelyn McCullough}
\date{September 2026}

\begin{document}

\maketitle

\section{Introduction}
\section{Background}
\subsection{State Sponsorship of Terrorism: Definition \& Context}

U.S. Congress's Gilmore Commission defined state-sponsored terrorism as ``the active involvement of a foreign government in training, arming, and providing other logistical and intelligence assistance as well as sanctuary to an otherwise autonomous terrorist group for the purpose of carrying out violent acts on behalf of the government against its enemies.'' \cite{gilmore1999assessing} When the US classifies a state as having sponsored a terrorist organization, it does so on the grounds that they ``repeatedly provided support for acts of international terrorism.'' \cite{usc22_2656f} Classifying a State as a state sponsor of terrorism allows the US to impose unilateral sanctions such as reduced foreign aid and a restriction on buying weapons. There are varying levels of active and passive state sponsorship.

Active sponsorship can be divided into three categories: control, coordination, and contact. Control sponsorship is where a group essentially performs acts on the state's behalf. Coordination sponsorship allows the groups to control their own interests and have some independence, but carry out acts that align with their sponsors' agendas. Contact sponsorship is where states are in regular communication with groups, but are not actively planning out attacks with them. This level of sponsorship is mainly for the purpose of maintaining a relationship for potential future cooperation. Active sponsorship decreased greatly between 1998 and 2000, which we can infer means that passive sponsorships increased.

Passive sponsorships contain three variations as well. The first is knowing toleration. This is where the government intentionally does not interfere with the group, but often means the country wants the group to carry out its goals without being directly connected to it. A very common form of sponsorship is unconcern, because countries do not see the group or organization as a threat. An example of this given by Byman is Canada's allowance of Liberation Tigers of Tamil Eelam to continue to raise funds without interfering. The final type of state sponsorship is incapacity. This is where a government does not have enough control over their country or region to control a groups' activities, therefore permitting it. \cite{article}

State sponsorship can be difficult to define because there are so many different ways it can be interpreted and executed. The binary way in which it will be described in this research cannot possibly capture the true complexities present in determining state sponsorship. In the dataset, a classification of 0 or 1 will be used to describe state sponsorship. This will lead to heterogeneity in our treatment group because any form of state sponsorship causes it to be included. So a state who provides active control sponsorship will be classified the same as a state providing passive incapacity sponsorship.

\subsection{Designation of Terrorist Organizations}
Before examining state sponsorship, it is important to consider how an organization is classified as a terrorist organization in the first place. Under Section 219 of the Immigration and Nationality Act, the U.S. Secretary of State may designate a Foreign Terrorist Organization (FTO) if it is foreign, engages in terrorist activity or retains the capability and intent to do so, and threatens U.S. nationals or national security. The State Department evaluates potential organizations using both classified and publicly available information. However, this designation process introduces an important consideration for my research: organizations that engage in political violence do not necessarily receive the same classification, particularly when their activities do not directly threaten U.S. interests. Furthermore, BAAD does not consist exclusively of U.S.-designated FTOs; its original sample was constructed using the Memorial Institute for the Prevention of Terrorism's database of terrorist organizations. Consequently, the comparison group in this study consists of organizations classified as nonsponsored within BAAD, rather than all violent organizations or all groups designated as terrorists by the U.S. These distinctions limit how broadly the findings can be generalized.

\subsection{Why States Sponsor Terrorist Organizations}
Governments often have several reasons for sponsoring a state organization. In \textit{Agents of Destruction}, they discuss how India supported the Liberation Tigers of Tamil Eelam (LTTE), which worked in favor of both parties. India delegated violence to this non-state actor because it allowed them plausible deniability, which allowed them to reduce the risk of direct conflict and preserve future negotiations between India and Sri Lanka. Meanwhile, LTTE received better equipment from India which greatly improved its prospects of gaining independence. 

\textit{Agents of Destruction} introduced the principal-agent framework with the State acting as the Principal and the Terrorist Organization acting as the Agent. The principal-agent problem in this context is that the state gives resources and support to the terrorist organization because it wants the organization to accomplish something, but the terrorist organization is an independent actor with its own goals. This leads to the states goals not always aligning with the terrorist organizations goals. This would help explain why some state/terrorist organization relationships are stronger than others. This is also shown with the LTTE because they were unhappy with the Indian government for working on a compromise between itself and Sri Lanka. This went directly against the groups' goals to achieve full independence. 

By being state sponsored, an organization has access to more resources and this therefore increases their capabilities, specifically for violence. The principal-agent relationship complicates this because, while the state is providing these resources, they do not have perfect control over how those resources are used. This relationship motivates the investigation into whether state-sponsored organizations exhibit different levels of lethality than non-sponsored organizations. 

\subsection{What Determines Terrorist Organization Lethality?}
In \textit{Dilettantes, Ideologues, and the Weak}, they determined that only 39\% of terrorist organizations between 1998 and 2005 killed anyone. This establishes the fact that just because a group is considered a 'terrorist organization', does not necessarily make them lethal. If most terrorist organizations are not killing people, then what distinguishes those that are. When a terrorist group is sponsored, they would have access to more resources, so could that possibly be the cause for the variations in lethality? 

In the paper, \textit{The Nature of the Beast}, they analyzed a variety of factors such as size, ideology, age, territorial control, alliance connections, and state sponsorship. Their model determined that size, ideology, territorial control and connectedness were important to lethality. A larger organization could be more lethal due to having more people to plan/execute, or have a greater variety of people willing to commit violence. A groups ideology could cause them to be more extremest and commit violent acts in the name of their mission. These are some already known characteristics that help explain why some groups are more lethal than others. 

In theory, sponsorship would provide more resources, leading to greater capabilities, however, prior research says that sponsorship was not an important predictor after accounting for other characteristics. When categorizing state sponsorship into a binary, it compresses both active and passive sponsorship into one category. This might mask a important distinction that different sponsor relationships could lead to different affects. When considering a principal-agent relationship, the state might be providing resources, but they do not get to necessarily control how they are being used. Because theory gives us a reason to expect sponsorship to matter, but previous empirical research complicates that explanation, the relationship between sponsorship and lethality warrants further investigation.

\subsection{Existing Evidence, Study Period, \& Research Question}
The time period of 1998 to 2005 crosses a threshold set by Collins in \textit{State-Sponsored Terrorism: In Decline, Yet Still a Potent Threat} between intermediate and low levels of state sponsorship. In 2001, he noted that state-sponsorship levels shifted from intermediate to low. This allows us to observe a period in which the broader nature of state sponsorship was changing.

Theory indicates there may be a causal relationship between state sponsorship and lethality. The binary classification of state sponsorship may be causing heterogeneity which could mask the affects under the principal agent problem. Prior research indicates that other characteristics matter and sponsorship has not previously emerged as a important predictor. Further investigation is worthwhile because there is a variation between theory and previous research, and discovering the underlying cause is important.

Specifically, the research question that will be asked is 'What is the effect of state sponsorship on the lethality of terrorist organizations from 1998 to 2005.' Theory leads to the expectation that state sponsorship increases the lethality of an organization. The analysis will determine whether the data provide evidence supporting this theoretical expectation. 

\section{Data}

\subsection{Data Sources, Unit of Observation, and Sample}

The primary source is the Big Allied and Dangerous Database I (BAAD I), which contains information on terrorist organizations and civilian fatalities during 1998--2005. The unit of observation is the terrorist organization, not an individual attack or an organization-year. The analytical sample consists of 395 organizations, including 32 classified by BAAD as state-sponsored (8.1\%) and 363 classified as nonsponsored (91.9\%). The outcome aggregates civilian fatalities for each organization from January 1, 1998, through December 31, 2005; it is not an annual fatality rate. All 395 organizations have an observed fatality count in the sample used here.

To investigate differences in forms of sponsorship, I also use monadic rebel-group active-support and de facto-support datasets. These supplemental records are organized by rebel group and year. I restrict both datasets to 1998--2005, merge them on the rebel-group identifiers and year, and summarize the resulting annual observations to the organization level so that they can be compared with BAAD. The restricted supplemental data contain 883 group-year records representing 160 distinct group identifiers. I match organizations by normalized names and verify additional documented aliases. Of BAAD's 32 sponsored organizations, 17 can be matched to the supplemental data. The remaining 15 are retained in the primary BAAD analysis but excluded from comparisons by type of support because their support classifications are unknown in the matched data. A missing match is not interpreted as evidence that an organization received no support.

\subsection{Variables and Measurement}

The dependent variable, \texttt{fatalities19982005}, is the total number of civilians killed by each organization during the study period, as defined in the BAAD I codebook. The primary explanatory variable, \texttt{statespond}, equals one when BAAD classifies an organization as supported by one or more governments and zero otherwise. This binary indicator does not distinguish among the types or intensity of government support discussed in the background section.

The supplemental data permit an exploratory distinction between active and de facto (referred to here as passive) support. For each organization, I count the observed years with at least one active sponsor (\texttt{NumNAG\_S} greater than zero) and the years with at least one de facto sponsor (\texttt{NumNAG\_DS} greater than zero). Organizations with active support but no recorded passive support are classified as ``Active only''; those with passive support but no recorded active support as ``Passive only''; and those with both types recorded at least once during 1998--2005 as ``Both.'' The two forms need not occur in the same year. These categories describe support observed in the supplemental datasets rather than replacing BAAD's original sponsorship indicator.

BAAD also provides organizational characteristics relevant to the relationship described in the background section: age in years as of December 31, 2005 (\texttt{OrgAge}); membership-size category (\texttt{ordsize}); an indicator of territorial control (\texttt{terrStrong}); and the number of alliance connections (\texttt{degree}). The size variable is categorical: code 0 includes groups with 0--100 members \emph{and} groups with uncertain membership estimates, while codes 1, 2, and 3 represent 100--1,000, 1,000--10,000, and at least 10,000 members, respectively. These variables will be considered in the subsequent identification strategy and analysis.

\subsection{Descriptive Statistics}

Table~\ref{tab:full_sample} summarizes the main outcome in the complete BAAD sample. Civilian fatalities are strongly right-skewed: the mean is 31.36, but the median is zero, and the largest recorded count is 3,505. The large standard deviation relative to the mean indicates substantial differences in lethality across organizations. This distribution is consistent with the background discussion that many organizations recorded no civilian fatalities during the study period.

\begin{table}[htbp]
\centering
\caption{Full-Sample Descriptive Statistics, 1998--2005}
\label{tab:full_sample}
\begin{tabular}{lr}
\hline
Statistic & Value \\
\hline
Organizations & 395 \\
State-sponsored organizations & 32 (8.1\%) \\
Nonsponsored organizations & 363 (91.9\%) \\
Mean civilian fatalities & 31.36 \\
Median civilian fatalities & 0 \\
Standard deviation & 202.04 \\
Minimum civilian fatalities & 0 \\
Maximum civilian fatalities & 3,505 \\
\hline
\end{tabular}
\end{table}

Table~\ref{tab:sponsorship_stats} compares the distribution of fatalities by BAAD's sponsorship classification. Sponsored organizations average 71.31 civilian fatalities, compared with 27.84 for nonsponsored organizations; the corresponding medians are 7.5 and zero. In terms of whether organizations recorded any civilian fatalities, 24 of the 32 sponsored organizations (75.0\%) recorded at least one fatality, compared with 131 of the 363 nonsponsored organizations (36.1\%). This distinction is important because a median of zero for nonsponsored organizations does not mean every organization in that group recorded zero fatalities. Nevertheless, the largest fatality count belongs to the nonsponsored category, which also has the higher standard deviation. These comparisons establish a descriptive relationship but do not show that sponsorship causes higher lethality.

\begin{table}[htbp]
\centering
\caption{Civilian Fatalities by BAAD State-Sponsorship Classification, 1998--2005}
\label{tab:sponsorship_stats}
\begin{tabular}{lrrrrr}
\hline
Sponsorship & $N$ & Mean & Median & SD & Maximum \\
\hline
Not sponsored & 363 & 27.84 & 0.0 & 207.21 & 3,505 \\
Sponsored & 32 & 71.31 & 7.5 & 125.04 & 458 \\
\hline
\end{tabular}
\end{table}

Figure~\ref{fig:sponsorship_boxplot} displays the comparison using the original recorded fatality counts. The high-fatality outliers make the boxes appear compressed near zero, especially for nonsponsored organizations, but the figure preserves the original scale of the outcome as requested. The log transformation is reserved for the regression analysis. This figure illustrates the unadjusted distribution and does not control for differences in organizational characteristics.

\begin{figure}[htbp]
    \centering
    \includegraphics[width=0.85\textwidth]{figure1_raw_fatalities.pdf}
    \caption{Civilian Fatalities by State Sponsorship, 1998--2005}
    \label{fig:sponsorship_boxplot}
    \par\smallskip
    {\small\textit{Note:} The vertical axis shows raw recorded civilian fatalities; extreme observations compress the boxes near zero.}
\end{figure}

\subsection{Exploratory Comparison of Support Types}

The distinction between active and passive sponsorship developed in the background section motivates a preliminary comparison within the 17 matched state-sponsored organizations. Table~\ref{tab:support_types} shows that 13 organizations have both types recorded, two have active support only, and two have passive support only. Organizations with both types have a mean of 136.85 fatalities and a median of 59; the active-only group has a mean of 5.50, and the passive-only group has a mean of 133.00. These descriptive differences are sensitive to individual organizations because the active-only and passive-only groups each contain only two observations. In addition, 15 of BAAD's 32 sponsored organizations could not be matched, so the table should not be interpreted as representative of all sponsored organizations.

\begin{table}[htbp]
\centering
\caption{Civilian Fatalities by Recorded Support Type: Matched Organizations, 1998--2005}
\label{tab:support_types}
\begin{tabular}{lrrrrrr}
\hline
Support type & $N$ & Mean & Median & SD & Minimum & Maximum \\
\hline
Active only & 2 & 5.50 & 5.5 & 6.36 & 1 & 10 \\
Both & 13 & 136.85 & 59.0 & 163.10 & 0 & 458 \\
Passive only & 2 & 133.00 & 133.0 & 188.09 & 0 & 266 \\
\hline
\end{tabular}
\end{table}

Taken together, these descriptive results motivate the subsequent analysis of whether the observed relationship between state sponsorship and civilian fatalities persists after accounting for organizational characteristics. The support-type comparison is exploratory because of its limited matched sample and is not used to infer causal effects.

\section{Identification Strategy}

\subsection{Research Question and Identification Problem}

The question I am trying to answer is whether state sponsorship affects the number of civilian fatalities attributed to terrorist organizations between 1998 and 2005. The treatment is BAAD's binary state-sponsorship classification, and the outcome is the total number of recorded civilian fatalities during that period. The main problem with identifying a causal effect is that organizations are not randomly selected to receive state support. States may be more likely to sponsor organizations that already have greater capabilities, larger networks, or goals that align with the state's interests. Those same characteristics could also affect how many fatalities an organization records. As a result, a difference in fatalities between sponsored and nonsponsored organizations does not necessarily mean that sponsorship caused the difference.

The data are organized at the organization level, so I cannot observe how an individual organization's lethality changed immediately before and after it received support. I also do not have a randomized treatment, a defensible instrumental variable, or a policy change that would support a difference-in-differences design. My strategy is therefore to compare organizations with observable similarities and examine how the sponsorship estimate changes across regression and matching specifications. This provides evidence about the association between sponsorship and lethality, but causal interpretation would still require an assumption that there are no relevant unobserved differences between the comparison groups.

\subsection{Regression Specifications}

I first estimate an unadjusted ordinary least squares regression:

\begin{equation}
\ln(1+\mathrm{Fatalities}_i)=\alpha+\beta\,\mathrm{Sponsored}_i+\varepsilon_i.
\end{equation}

I use $\ln(1+\mathrm{Fatalities})$ because the raw fatality counts are highly right-skewed and many organizations recorded zero fatalities. Adding one allows those organizations to remain in the regression. In the baseline model, $\beta$ measures the difference in average log-transformed fatalities between sponsored and nonsponsored organizations, without accounting for other characteristics.

I then estimate models that add organizational age and ideology, organizational capacity, and finally all of the selected controls:

\begin{align}
\ln(1+\mathrm{Fatalities}_i) ={}& \alpha+\beta\,\mathrm{Sponsored}_i
+\gamma_1\mathrm{OrgAge}_i
+\gamma_2\mathrm{Religious}_i \\
&+\gamma_3\mathrm{Ethnonationalist}_i
+\sum_{k=1}^{3}\delta_k\mathbf{1}(\mathrm{Size}_i=k)
+\theta\mathrm{Territorial}_i
+\lambda\mathrm{Degree}_i+\varepsilon_i.
\end{align}

Organizational size is treated as a categorical variable, with size category zero as the reference group. I use HC3 heteroskedasticity-robust standard errors because fatality outcomes vary substantially across organizations. All four OLS specifications use the same 395 organizations, so changes in coefficients are not caused by observations being dropped for missing values.

There is an important limitation to the fully adjusted specification. Size, territorial strength, and alliance connections could be factors that existed before sponsorship, but they could also be affected by sponsorship itself. If state support helps an organization grow or expand its network, controlling for those characteristics could remove part of the relationship I am trying to measure. For this reason, I compare specifications with and without these variables rather than treating the fully adjusted coefficient as a definitive causal effect.

\subsection{Matching Strategy and Assumptions}

As a second approach, I use propensity score matching to construct a comparison group of nonsponsored organizations that are more similar to the sponsored organizations in observed characteristics. My primary matching specification estimates the probability of sponsorship using organizational age, religious ideology, and ethnonationalist ideology. I use one-to-one nearest-neighbor matching without replacement, an ATT-oriented design, and a caliper of 0.2 standardized units of the estimated propensity score. I then calculate fatality differences within each sponsored--nonsponsored pair. This design is intended to estimate an average difference for the sponsored organizations that could be matched; it does not automatically identify the effect for every sponsored organization.

For matching to have a causal interpretation, organizations would need to be comparable after conditioning on the chosen characteristics, and sponsored organizations would need suitable nonsponsored comparisons. The first assumption is difficult to defend because the dataset does not measure every reason that states choose particular organizations to support. In addition, even age and ideology need to be interpreted carefully because their timing relative to the start of sponsorship is not established here. I evaluate overlap and covariate balance rather than assuming that matching removes these problems. A secondary matching specification also includes organizational capacity, but I treat that as a sensitivity analysis because those variables may be affected by sponsorship.

\section{Results}

\subsection{Initial OLS Analysis}

The initial regression shows a positive association between state sponsorship and log-transformed fatalities. In the unadjusted specification, the sponsorship coefficient is 1.645 and statistically significant using HC3 robust standard errors. However, the estimate becomes smaller when I account for organizational characteristics. Table~\ref{tab:part3_ols} reports the sponsorship coefficient from each of the four OLS specifications.

\begin{table}[htbp]
\centering
\caption{State Sponsorship Coefficients Across OLS Specifications}
\label{tab:part3_ols}
\begin{tabular}{lrrrr}
\hline
Specification & Estimate & HC3 SE & $p$-value & $N$ \\
\hline
Unadjusted & 1.645 & 0.387 & $<0.001$ & 395 \\
Organizational controls & 0.289 & 0.358 & 0.420 & 395 \\
Age and ideology controls & 0.903 & 0.413 & 0.029 & 395 \\
Full controls & 0.089 & 0.360 & 0.805 & 395 \\
\hline
\end{tabular}
\par\smallskip
{\small\textit{Note:} The dependent variable is $\ln(1+\text{recorded fatalities})$. Organizational controls include age, categorical size, territorial strength, and network degree. The age and ideology model includes age, religious ideology, and ethnonationalist ideology. The full model includes all listed controls. Standard errors are HC3 heteroskedasticity-robust.}
\end{table}

The coefficient decreases to 0.289 when organizational capacity characteristics are included and to 0.089 in the fully adjusted model, where it is no longer statistically significant ($p=0.805$). In contrast, sponsorship remains statistically significant in the model that adjusts only for age and ideology ($\hat\beta=0.903$, $p=0.029$). The fully adjusted model explains approximately 45.0\% of the variation in the log-transformed outcome, compared with 7.5\% for the unadjusted model. Size, network degree, territorial strength, and religious ideology have positive and statistically significant coefficients in the fully adjusted model. The differences across these specifications suggest that organizational characteristics are closely related to the observed sponsorship association, although they do not tell us whether those characteristics are confounders or consequences of sponsorship.

\subsection{Initial Matching Analysis}

The primary matching procedure retained 30 of the 32 sponsored organizations and selected 30 nonsponsored controls. Balance improved substantially for organizational age and religious ideology, but ethnonationalist ideology remained imbalanced. This means the matched organizations are more similar on some observed characteristics, but they are not completely balanced. Table~\ref{tab:part3_matching} summarizes the matched differences on the original fatality scale.

\begin{table}[htbp]
\centering
\caption{Matched Differences in Recorded Fatalities}
\label{tab:part3_matching}
\begin{tabular}{lrrrr}
\hline
Matching specification & Pairs & Mean difference & 95\% CI & $p$-value \\
\hline
Age and ideology & 30 & 25.67 & $[-28.07,\ 79.41]$ & 0.337 \\
Capacity-adjusted & 26 & -122.23 & $[-406.76,\ 162.29]$ & 0.385 \\
\hline
\end{tabular}
\par\smallskip
{\small\textit{Note:} Differences are sponsored minus nonsponsored fatalities within matched pairs. Confidence intervals are based on pair-level $t$ tests. The secondary specification includes potentially post-treatment organizational capacity variables and retains a different subset of organizations.}
\end{table}

The primary matched estimate is positive, with sponsored organizations recording an average of 25.67 more fatalities than their matched controls, but its 95\% confidence interval includes zero ($p=0.337$). The alternative matching specification, which includes capacity characteristics, produces a negative estimate that is also statistically insignificant. Because the two approaches retain different sets of organizations, the change in sign cannot be attributed only to the additional controls. The matching results therefore show uncertainty about both the direction and magnitude of the sponsorship association.

\subsection{Interpretation and Limitations of the Initial Findings}

At this stage, the results do not provide a consistent statistically significant relationship between sponsorship and lethality after accounting for observable organizational differences. The unadjusted comparison is positive, but the adjusted regression and primary matched comparison are not statistically significant at the 5\% level. This does not prove that state sponsorship has no effect. The sponsorship measure groups several kinds of support into one binary indicator, the support-type data cover only 17 of the 32 sponsored BAAD organizations, and matching cannot account for unobserved factors that may influence both sponsorship and violence. These limitations are especially important for the identification question. My initial analysis establishes how the observed relationship changes with model specification, but the available data do not allow me to isolate a definitive causal effect.


\end{document}
